Lightweight image super-resolution must recover long-range repeated structure under strict parameter and latency budgets. Existing content-aware grouping methods improve non-local interaction, but they still process each subgroup with a largely homogeneous attention pattern and therefore waste computation on weak correspondences. We propose \method{}, a lightweight SR network built around three coordinated components: a Dynamic Prototype Router (\dpr{}) for image-adaptive content organization, Progressive Focused Sparse Attention (\pfsa{}) for subgroup-dependent connection filtering, and Low-to-High Multi-Level Reconstruction (\lmr{}) for coarse-to-fine restoration. The resulting \moduleName{} first stabilizes large-scale structure, then refines medium-scale transitions, and finally restores high-frequency details with inherited attention priors. Experiments on Set5, Set14, BSD100, Urban100, and Manga109 show that \method{} is especially effective on complex benchmarks. Under \(\times4\) SR, it achieves 27.24 dB on Urban100 and 31.55 dB on Manga109 with only 567K parameters and 46.1G FLOPs, improving \baseline{} by 0.37 dB and 0.24 dB on the two datasets while slightly reducing FLOPs and latency. These results show that content-aware grouping benefits substantially from scale-aware focused interaction rather than a single dense subgroup attention stage.
